\section{Background and Related Work}
\label{section_background}

\subsection{Temporal Distribution Shift}
A central challenge in real-world machine learning systems is that the data-generating process is rarely stationary. As models are deployed over months, years, or decades, both inputs and label semantics may evolve, causing systematic discrepancies between distributions encountered during training and those observed at inference. Such temporal evolution, commonly referred to as \emph{temporal distribution shift}, is pervasive across domains. Understanding the structure of this shift is essential for assessing and improving temporal robustness. However, prior work has largely focused on static or domain-level distribution shifts \cite{ben2010theory,gulrajani2021in_search_of_lost,koh2021wilds}, with limited attention to drift that unfolds sequentially over time, particularly in non-generated, in-the-wild settings \cite{yao2022wildtime}.

Following established terminology in concept drift research \cite{gama2014a_survey_on_concept_drift}, temporal distribution shift can be characterized along three complementary axes (Figure \ref{fig:distribution_shift}):

\textbf{Covariate shift} occurs when the input distribution $P(X)$ changes while the conditional $P(Y \mid X)$ remains fixed. This type of shift commonly arises in natural data streams where observational setups, sociocultural conventions, or user behavior gradually evolve.

\textbf{Label shift} refers to changes in the marginal label distribution $P(Y)$. Long-term textual or behavioral datasets frequently exhibit such imbalance drift as the prevalence of topics, categories, or rating patterns changes over time.

\textbf{Concept drift} denotes changes in the conditional distribution $P(Y \mid X)$, implying that identical inputs may correspond to different labels at different time points. This form of drift is particularly pronounced in domains where semantics or visual attributes evolve, such as historical portraits \cite{ginosar2015century} or online platforms.

In practice, these forms of drift rarely occur in isolation and often manifest in combination \cite{aguiar2023comprehensiveanalysisconceptdrift}. As a result, the temporal robustness of a model depends not only on the magnitude of drift but also on how its architectural assumptions and inductive biases interact with the evolving data distribution.

\subsection{Model Robustness to Distribution Shift}
% Shared by paper/main-lncs.tex and paper/main-preprint.tex.
% Define a command to draw stars
\newcommand{\drawstar}[2]{
    \draw[fill=#1] (#2) -- ++(0.1,0.03) -- ++(0.03,0.1) -- ++(0.03,-0.1) -- ++(0.1,-0.03) -- ++(-0.1,-0.03) -- ++(-0.03,-0.1) -- ++(-0.03,0.1) -- cycle;
}
\begin{figure}[t]
    \centering
    \resizebox{\linewidth}{!}{
        \begin{tikzpicture}
            % Axes
            \draw[->] (-0.2,0) -- (1.5,0) node[below] {};
            \draw[->] (0,-0.2) -- (0,1.5) node[left] {};

            % Decision boundary (black line)
            \draw[thick] (0.1, 0.25) .. controls (0.5, 0.2) and (0.9, 1)  .. (1.4, 1);

            % Blue dots (class 1) with black border
            \filldraw[black, fill=accent_blue] (0.3, 0.5) circle (0.1);
            \filldraw[black, fill=accent_blue] (0.65, 1) circle (0.1);
            \filldraw[black, fill=accent_blue] (1.2, 1.2) circle (0.1);

            % Red stars (class 2)
            \drawstar{accent_orange}{0.45, 0.25};
            \drawstar{accent_orange}{0.85, 0.5};
            \drawstar{accent_orange}{1.05, 0.25};
            \drawstar{accent_orange}{1.25, 0.75};

            % Label
            \node at (0.75, 1.7) {\scriptsize Original};
        \end{tikzpicture}

        \hspace{0.15cm}

        \begin{tikzpicture}
            % Concept Drift
            % Axes
            \draw[->] (-0.2,0) -- (1.5,0) node[below] {};
            \draw[->] (0,-0.2) -- (0,1.5) node[left] {};

            % Decision boundary (black line)
            \draw[thick] (0.1, 0.7) .. controls (0.9, 1) and (1.3, 0.5)  .. (1.4, 0.42);

            % Blue dots (class 1) with black border
            \filldraw[black, fill=accent_blue] (0.65, 1) circle (0.1);
            \filldraw[black, fill=accent_blue] (1.2, 1.2) circle (0.1);
            \filldraw[black, fill=accent_blue] (1.25, 0.75) circle (0.1);

            % Red stars (class 2)
            \drawstar{accent_orange}{0.3, 0.5};
            \drawstar{accent_orange}{0.45, 0.25};
            \drawstar{accent_orange}{0.85, 0.5};
            \drawstar{accent_orange}{1.05, 0.25};

            % Label
            \node at (0.75, 1.7) {\scriptsize Concept Drift};
        \end{tikzpicture}

        \hspace{0.15cm}

        \begin{tikzpicture}
            % Virtual Drift
            % Axes
            \draw[->] (-0.2,0) -- (1.5,0) node[below] {};
            \draw[->] (0,-0.2) -- (0,1.5) node[left] {};

            % Decision boundary (black line)
            \draw[thick] (0.1, 0.25) .. controls (0.5, 0.2) and (0.9, 1)  .. (1.4, 1);

            % slightly shifted
            % \draw[thick] (0.1, 0.15) .. controls (0.8, 0.2) and (0.9, 1)  .. (1.4, 1);

            % Blue dots (class 1) with black border
            \filldraw[black, fill=accent_blue] (0.3, 0.55) circle (0.1);
            \filldraw[black, fill=accent_blue] (0.25, 0.75) circle (0.1);
            \filldraw[black, fill=accent_blue] (0.5, 0.85) circle (0.1);

            % Red stars (class 2)
            \drawstar{accent_orange}{1, 0.5};
            \drawstar{accent_orange}{1.25, 0.35};
            \drawstar{accent_orange}{1.3, 0.7};

            % Label
            \node at (0.75, 1.7) {\scriptsize Covariate Drift};
            \node at (0.75, 1.4) {\tiny (Virtual Drift)};
        \end{tikzpicture}

        \hspace{0.15cm}

        \begin{tikzpicture}
            % Prior Probability Shift
            % Axes
            \draw[->] (-0.2,0) -- (1.5,0) node[below] {};
            \draw[->] (0,-0.2) -- (0,1.5) node[left] {};

            % Decision boundary (black line)
            \draw[thick] (0.1, 0.25) .. controls (0.5, 0.2) and (0.9, 1)  .. (1.4, 1);

            % slightly shifted
            % \draw[thick] (0.2, 0.15) .. controls (0.4, 0.1) and (0.8, 1)  .. (1.4, 0.9);

            % Blue dots (class 1) with black border
            \filldraw[black, fill=accent_blue] (0.3, 0.5) circle (0.1);
            \filldraw[black, fill=accent_blue] (0.65, 1) circle (0.1);
            \filldraw[black, fill=accent_blue] (1.2, 1.2) circle (0.1);
            \filldraw[black, fill=accent_blue] (0.35, 1.1) circle (0.1);
            \filldraw[black, fill=accent_blue] (0.4, 0.8) circle (0.1);
            \filldraw[black, fill=accent_blue] (0.6, 1.2) circle (0.1);
            \filldraw[black, fill=accent_blue] (0.15, 0.7) circle (0.1);
            \filldraw[black, fill=accent_blue] (0.25, 0.85) circle (0.1);

            % Red stars (class 2)
            \drawstar{accent_orange}{0.85, 0.5};
            \drawstar{accent_orange}{1.05, 0.25};

            % Label
            \node at (0.75, 1.7) {\scriptsize Label Drift};
        \end{tikzpicture}

    }
    \caption{Drift types in a binary classification setting. Circles and stars indicate the label classes; the curve represents the decision boundary. \textit{Concept drift} induces a change in the decision boundary, \textit{Covariate drift} (virtual drift) changes the input distribution, and \textit{Label drift} alters the relative class frequency \iflncs\cite{gama2014a_survey_on_concept_drift}\else\cite{holzinger2024analysis}\fi.}
    \tightvspace{-5pt}
    \label{fig:distribution_shift}
    \tightvspace{-0.2cm}
\end{figure}


Understanding how learning algorithms behave under distribution shift has become an important research direction. Early work approached robustness through domain adaptation \cite{ben2010theory} and out-of-distribution generalization \cite{gulrajani2021in_search_of_lost}, formalizing shift as a transition between a small number of discrete source and target domains. The introduction of large-scale benchmarks such as \texttt{WILDS} \cite{koh2021wilds}, which focuses on naturally occurring distribution shifts without explicit temporal indexing, and more recently \texttt{Wild-Time} \cite{yao2022wildtime}, which explicitly models time-indexed data and temporal distribution shifts, has broadened this perspective by providing evaluation protocols that more closely reflect real-world deployment scenarios.

A complementary line of research has examined robustness through model diagnostics and monitoring. Methods for detecting drift onset by examining changes in latent representations or predictive uncertainty have shown promise in deep learning settings \iflncs\cite{rabanser_failing_2019}\else\cite{rabanser_failing_2019,ackerman2021automaticallydetectingdatadrift}\fi. Recent empirical analyses have characterized how concept drift manifests in practice, including its locality and temporal progression across large-scale data streams \cite{aguiar2023comprehensiveanalysisconceptdrift}. At the systems level, work on data-centric and continual-learning infrastructures has explored how to maintain model quality over time through cost-aware retraining and pipeline orchestration \iflncs\cite{mahadevan_cost_aware_2024,boether_modyn_2025}\else\cite{mahadevan_cost_effective_2023,mahadevan_cost_aware_2024,tian_continuum_2018,boether_modyn_2025,holzinger2024analysis}\fi.

Despite these advances, robustness studies commonly focus either on a single architecture evaluated across multiple datasets or on a single dataset used to compare a narrow set of architectures. As a result, comparatively little is known about how architectural design choices interact with long-range temporal drift across heterogeneous modalities and tasks. This gap is particularly relevant given the diversity of inductive biases exhibited by modern neural models: convolutional networks encode locality and translation equivariance, recurrent networks capture sequential structure \cite{hochreiter1997long,cho2014learning}, and Transformer-based encoders rely on self-attention with minimal structural priors \cite{vaswani2017attention}. Likewise, large-scale pretraining in vision and language \cite{devlin2019bert,liu2019roberta} introduces representations shaped by broad historical data, yet their behavior under extended temporal drift remains poorly characterized.

The present study addresses this open question by comparing these architectural families under a shared temporal evaluation protocol and across multiple modalities, enabling a controlled analysis of how model design influences robustness under real-world temporal distribution shift.

\subsection{Robustness Evaluation vs.\ Temporal Adaptation}

This work evaluates the inherent temporal robustness of neural architectures under distribution shift, measuring how different model families, trained only on cumulative historical data, perform as the temporal gap between training and evaluation widens. This isolates the robustness arising from architectural design and pretraining alone, prior to any adaptive intervention.

A complementary line of research instead investigates how models can \emph{adapt} once drift is detected, \iflncs through scheduled or cost-aware retraining \cite{mahadevan_cost_aware_2024} and continual-learning pipelines \cite{boether_modyn_2025}\else through scheduled or cost-aware retraining \cite{mahadevan_cost_effective_2023,mahadevan_cost_aware_2024}, continual-learning pipelines \cite{tian_continuum_2018,boether_modyn_2025,holzinger2024analysis}, or accuracy-aware data maintenance \cite{ralf2023wooders}\fi, addressing when to retrain, how much data to incorporate, and how to trade performance against cost. Our study is orthogonal to that literature and provides a foundation on which such adaptation strategies can be designed, evaluated, and compared.
